\documentclass[12pt,a4paper]{article}
\usepackage[utf8]{inputenc}
\usepackage[spanish]{babel}
\usepackage{amsmath}
\usepackage{amsfonts}
\usepackage{amssymb}
\usepackage{graphicx}
\usepackage{hyperref}
\usepackage{geometry}
\geometry{left=2.5cm,right=2.5cm,top=2.5cm,bottom=2.5cm}

\title{Proyecto LINX: Diseño Conceptual de UAV Solar de Gran Altitud y Larga Autonomía (HALE)}
\author{Josua Emiliano Alamilla Jimenez \and Roció Nayeli Oliva González \and Asesor: Dr. Gustavo Medina Tanco}
\date{\today}

\begin{document}

\maketitle

\begin{abstract}
Este documento sintetiza el alcance, la metodología y el progreso actual del proyecto LINX, el cual tiene como objetivo el diseño de un Vehículo Aéreo No Tripulado (UAV) de tipo HALE (High Altitude and Long Endurance) alimentado por energía solar, destinado a operar en la estratosfera de manera autónoma y autosuficiente.
\end{abstract}

\section{Alcance del Proyecto (Scope)}
El objetivo principal del proyecto LINX es diseñar, desarrollar y validar un UAV propulsado por energía solar capaz de transportar y mantener operativa una carga útil de hasta 5 kg a gran altitud (estratosfera) durante periodos prolongados (meses).
\begin{itemize}
    \item \textbf{Autonomía:} Totalmente autónomo con capacidad de navegación, estabilización y gestión energética.
    \item \textbf{Región de Operación:} Cercano a la latitud de México, entre 19$^\circ$ y 24$^\circ$ N.
    \item \textbf{Especificaciones Clave:} Peso máximo de despegue (MTOW) de 30 kg, envergadura de 13.384 m, y superficie alar de 14.928 m$^2$.
    \item \textbf{Marco Legal:} Operación bajo NOM-107-SCT3-2019, requiriendo permisos para operación BVLOS y autorización para Aeronave Experimental por encima de 400 ft.
\end{itemize}

\section{Metodología y Dimensionamiento}
El proyecto se divide en 5 fases, avanzando desde el diseño conceptual hasta la manufactura y pruebas:
\begin{enumerate}
    \item \textbf{Fase A:} Finalizada.
    \item \textbf{Fase B (Diseño Conceptual):} Dimensionamiento inicial de pesos (carga útil, estructura, baterías, etc.), selección de configuración y análisis de viabilidad (38\% del esfuerzo).
    \item \textbf{Fase C (Diseño Detallado):} Refinamiento aerodinámico y estructural, selección de componentes y simulaciones detalladas (19\% del esfuerzo).
    \item \textbf{Fase D (Manufactura y Pruebas):} Construcción de prototipos y modelos de ingeniería, pruebas en tierra y vuelo (38\% del esfuerzo).
    \item \textbf{Fase E (Documentación y Entrega):} Reporte final (6\% del esfuerzo).
\end{enumerate}

\subsection{Ecuaciones de Dimensionamiento y Rendimiento}
Partiendo de la información estructurada, el diseño inicial estima el peso total de despegue ($W_{TO}$) como la suma de sus componentes:
\begin{equation}
    W_{TO} = W_{PL} + W_{S} + W_{B} + W_{E}
\end{equation}
Donde $W_{PL}$ es el peso de carga útil, $W_S$ es el peso de sistemas, $W_B$ es el peso de baterías y $W_E$ es el peso vacío.

Para el análisis aerodinámico y el equilibrio en vuelo, la fuerza de sustentación ($L$) se iguala al peso:
\begin{equation}
    L = W = \frac{1}{2} \rho v^2 S C_L
\end{equation}

El cálculo de la potencia requerida y disponible es fundamental. La potencia en el eje (Brake Horse Power, BHP) para el motor eléctrico se estima como:
\begin{equation}
    BHP = \frac{V \cdot I \cdot \eta_{motor}}{746}
\end{equation}

El empuje ($T$) proporcionado por las hélices y la eficiencia propulsiva ($\eta_P$) se calculan mediante:
\begin{equation}
    T = C_T \rho n^2 D^4
\end{equation}
\begin{equation}
    \eta_P = \frac{T \cdot V}{BHP}
\end{equation}

Además, el balance energético (Power Budget) toma en consideración las eficiencias de la cadena de potencia ($\eta_{cadena} = \eta_{prop} \cdot \eta_{motor} \cdot \eta_{ESC}$).

Para el desarrollo aerodinámico se ha optado por el método de paneles (XFLR5) y CFD. 

\section{Toma de Decisiones y Análisis de Riesgos}

\subsection{Descomposición Funcional y QFD}
El proceso de diseño partió de un \textbf{Análisis de Descomposición Funcional} exhaustivo, mapeando la función global a requerimientos específicos de los subsistemas. Para traducir estos requerimientos y priorizarlos, se emplearon herramientas como el Despliegue de la Función de Calidad (QFD) en conjunto con el \textbf{Proceso de Jerarquía Analítica (AHP de Saaty)}. 

Se priorizaron dos grandes rubros:
\begin{itemize}
    \item \textbf{Aerodinámica:} Los requerimientos con mayor peso fueron generar sustentación suficiente en crucero (R04, peso 0.0704) y reducir el arrastre parásito para maximizar la relación L/D (R05, peso 0.0704).
    \item \textbf{Manufacturabilidad y Estructura:} Se dio prioridad a limitar las deformaciones para mantener la integridad aerodinámica (R19, peso 0.0405) y a la validación estructural ante cargas y vibraciones (R22, peso 0.0386).
\end{itemize}

\subsection{Análisis de Modos de Falla y Efectos (FMEA)}
Para garantizar la fiabilidad del UAV en misiones BVLOS y de larga duración, se implementó la metodología \textbf{FMEA}. Se evaluó el Número de Prioridad de Riesgo (RPN) multiplicando la Probabilidad (P), Severidad (S) y Detección (D). 
Los riesgos más críticos identificados fueron:
\begin{enumerate}
    \item \textbf{Déficit energético nocturno (RPN: 15):} Crítico durante el crucero estratosférico, mitigado mediante un estricto balance de batería y eficiencia solar.
    \item \textbf{Fallo de Motor/ESC (RPN: 12):} Riesgo mayor durante la fase de ascenso.
    \item \textbf{Pérdida de Enlace C2 (RPN: 12):} Asociado a las misiones de larga distancia y gran altitud.
    \item \textbf{Aterrizaje Forzoso / Hard Landing (RPN: 12):} Debido a las dimensiones del ala y la ligereza estructural.
\end{enumerate}

\section{Trabajo Realizado e Integración}
Hasta la fecha, se ha avanzado significativamente en el diseño conceptual y detallado:
\begin{itemize}
    \item \textbf{Configuración Aerodinámica:} Se evaluó la configuración de fuselaje integrado con estabilizador y configuración "Tandem Pusher". Se definieron perfiles alares y se optimizó el ángulo diedro y alabeo a lo largo del día para mejorar la captación solar.
    \item \textbf{Sistema Propulsivo:} Se seleccionó una hélice con diámetro de 1.35 m, coeficiente de empuje $C_T=0.1411$, operando a 50 rev/s, con una potencia real por motor de 6720 W.
    \item \textbf{Sistema de Potencia (Solar):} Se realizó un benchmarking exhaustivo de celdas solares limitadas a 5 kg. Las celdas Spectrolab fueron elegidas por ofrecer el mayor margen de diseño y cumplir los voltajes del bus, considerando las horas solares efectivas (solsticio de invierno) en la latitud 19.4$^\circ$.
    \item \textbf{Estado del Arte y Referencias:} Se cuenta con una amplia recopilación de datos, incluyendo libros de diseño (Raymer, Aircraft Design), artículos sobre UAVs solares y análisis de estabilidad (revisados desde el repositorio bibliográfico del proyecto).
\end{itemize}

\section{Conclusión}
El proyecto LINX representa un avance importante en la aeronáutica nacional, integrando sistemas solares, diseño estructural liviano y metodologías de ingeniería de sistemas rigurosas para lograr un vuelo estratosférico sostenido.

\end{document}
